% ====================================================================
% graphite_ica_consolidated_ch1_G1.tex  —  합류(Consolidation) Chapter 1 (G1)
%   트렁크 = (B) 전하보존 6장 계열, 엔진 이식 = (A) relaxation-length spectrum.
%   keystone: χA = Ch1 Level A(mobility 가속) / 방향성 barrier lowering 은 Ch3 Level B.
%   지배 표준: 계획서 §00 GS-1~4 (학부 무비약 / 전문 깊이 / grounded 가정 / 과도가정 배제).
%   출처: PHASE_DIAG_INTENT_GAP_RESULT, _SALVAGE_LEDGER, _REFS67_DOSSIER,
%         2026-05-29-consolidation-roadmap.  (기존 (A)(B) 원본은 보존, 본 파일은 신규 G1.)
% Date: 2026-05-29   Author: Project_Anode_Fit (Claude 측)
% ====================================================================
\documentclass[11pt,a4paper]{article}
\usepackage[margin=22mm]{geometry}
\usepackage{kotex}
\usepackage{amsmath,amssymb,mathtools,bm}
\usepackage{booktabs,longtable,array,tabularx}
\usepackage{enumitem}
\usepackage{hyperref}
\usepackage{amsthm}
\usepackage{mdframed}
\usepackage{fancyhdr}
\usepackage{lastpage}
\usepackage{setspace}
\setstretch{1.12}
\setlength{\parskip}{0.4em}
\setlength{\parindent}{0pt}
\setlength{\headheight}{14pt}
\hypersetup{colorlinks=true, linkcolor=blue, urlcolor=blue, citecolor=blue,
  pdftitle={Graphite ICA tail — 합류 Chapter 1 (G1)}}
\pagestyle{fancy}\fancyhf{}
\lhead{흑연 음극 ICA tail — 합류 Chapter 1 (G1)}\rhead{\thepage/\pageref{LastPage}}
\renewcommand{\headrulewidth}{0.4pt}

\newtheorem*{groundingbox}{Grounding}
\newtheorem*{boundbox}{유효범위(Validity bound)}
\newtheorem*{flagbox}{FLAGGED (미확정)}
\newtheorem*{keybox}{Keystone}

\newcommand{\dd}{\mathrm{d}}
\newcommand{\eff}{\mathrm{eff}}
\newcommand{\eq}{\mathrm{eq}}
\newcommand{\app}{\mathrm{app}}
\newcommand{\drive}{\mathrm{drive}}
\newcommand{\tot}{\mathrm{tot}}
\newcommand{\bg}{\mathrm{bg}}
\newcommand{\cell}{\mathrm{cell}}
\newcommand{\OCV}{\mathrm{OCV}}
\newcommand{\simplecase}{\mathrm{simple}}
\newcommand{\AL}[1]{\textsf{[AL-#1]}}

\title{\textbf{리튬이온전지 흑연 음극 ICA($dQ/dV$) tail 의 열역학 무대 $\cdot$ 동역학 주역 이론}\\[0.3em]
\large 합류 Chapter 1 (G1): 전하보존 $V_n$ + 유효 mobility + relaxation-length spectrum + 피팅식}
\author{Project\_Anode\_Fit (Claude 측)}
\date{2026-05-29}

\begin{document}
\maketitle

% ====================================================================
\section*{서(序) — 관찰에서 출발하는 단일 인과 사슬 (N-1)}
\addcontentsline{toc}{section}{서 — 단일 인과 사슬}
% ====================================================================
본 장은 다음 \textbf{하나의 관찰}에서 출발하여 끝까지 그 끈을 놓지 않는다.

\begin{quote}
\emph{흑연 음극의 $dQ/dV$(ICA) 곡선에서, 상변이 peak 의 꼬리가 \textbf{온도가 낮을수록 길게
늘어져 다음 peak 와 겹치고}, 높을수록 짧게 끝난다. 또 전류(C-rate)가 커지면 peak 겹침이
심해진다.}
\end{quote}

이 관찰을 다섯 단계로 푼다 (각 단계는 학부 수준으로 따라갈 수 있게 무비약으로 전개; 본 장의
지배 표준 GS-1).
\begin{enumerate}[label=(\arabic*)]
\item \textbf{평형 열역학만으로는 안 된다.} 평형 전이폭은 $w=RT/F$ 라 \emph{저온서 오히려
  좁아져}(\S\ref{sec:eq}), 평형 이론은 ``저온 $\to$ 짧은 꼬리''를 예측한다. 관찰과 \emph{반대}.
\item \textbf{따라서 꼬리는 동역학이다.} 진행이 평형을 \emph{유한 속도}로 따라가며 생기는 lag 가
  꼬리를 만든다. 속도상수 $k=k_0e^{-\Delta G_a/RT}$ 는 저온서 작아져 lag 가 커지고 꼬리가
  길어진다(\S\ref{sec:rate}). 관찰과 \emph{일치}.
\item \textbf{그 배리어는 극판 전위로 바뀐다.} 전류가 흐를 때 내부 전위는 OCV 평형과 다르며,
  이 전위 구동력이 \emph{진행 속도(mobility)}를 바꾼다(\S\ref{sec:barrier}). 그래서 C-rate 가
  peak 겹침을 키운다.
\item \textbf{무대는 전하보존이 깐다.} 내부 음극 전위 $V_n$ 은 외부에서 읽는 값이 아니라
  \emph{전하 보존식의 해}로 결정된다(\S\ref{sec:charge}). 꼬리의 \emph{깊이}는 실제 전극의
  barrier 분포가 만드는 \emph{relaxation-length spectrum}으로 설명된다(\S\ref{sec:spectrum}).
\item \textbf{결과 = 피팅 가능한 닫힌 논리식}(\S\ref{sec:fiteq}, 식~\eqref{eq:fiteq}). 이후
  발열(Ch2,4)$\cdot$히스테리시스(Ch5)$\cdot$수치검증(Ch6)으로 확장한다.
\end{enumerate}

\noindent\textbf{한 문장 요지: 열역학이 \emph{무대}($V_n$, 평형 target)를 깔고, 관찰된 꼬리
거동의 \emph{주역}은 동역학이다.}

\begin{groundingbox}
\textbf{지배 표준(GS).} (GS-1) 모든 식은 물리적 의미를 prose 로 먼저 제시한 뒤 수식화하고
중간 단계를 생략하지 않는다. (GS-2) 결론은 출판 수준 엄밀성. (GS-3) 모든 가정은 Assumption
Ledger(AL-\#: 논문/교과서 근거)를 인용; 근거 없는 항목은 FLAGGED. (GS-4) 임의 clip/cap/
softplus/step-function 정의식 금지; 속도 제한은 물리적 병목으로, BOUNDED 가정은 유효범위 병기.
\end{groundingbox}

\tableofcontents
\newpage

% ====================================================================
\section{명명 $\cdot$ 기호 규약 (G0)}\label{sec:notation}
% ====================================================================
\textbf{역사적 명칭 매핑(P3-7).} 본 합류본은 기존 두 계열을 통합한다: (A) relaxation-length
spectrum kernel-integral 계열(기호 $\theta$), (B) 전하보존 6장 계열(기호 $\xi$). 본 G1 은
\textbf{트렁크 (B)의 진행률 기호 $\xi$ 로 통일}한다; (A)의 $\theta$ 는 본 장의 $\xi$ 와 동일물.
``$ver.1$--$ver.5$''(역사적 파일 이력)와 ``Chapter 1--6''(본 구조)을 혼동하지 않는다.

\begin{longtable}{p{0.15\textwidth} p{0.12\textwidth} p{0.65\textwidth}}
\toprule 기호 & 단위 & 의미 \\ \midrule \endhead
$q$ & --- & 방전 진행 좌표 $q=Q_\mathrm{ext}/Q_\cell$ \\
$Q_\mathrm{ext},Q_\cell$ & C & 외부 누적 charge; 기준 cell 용량. $|I|=dQ_\mathrm{ext}/dt$ \\
$V_n$ & V & 내부 음극 전위 — \textbf{전하 보존식의 root}(외부 lookup 아님) \\
$V_{n,\app}$ & V & apparent(측정) 전위 $=V_n+s_I|I|R_n$ \\
$V_{n,\drive}$ & V & 속도식 구동 전위; 기본 $=V_n$ \\
$V_{n,\OCV}$ & V & 평형($\xi_j=\xi_{\eq,j}$)에서의 charge-balance 해 \\
$\xi_j,\ \xi_{\eq,j}$ & --- & $j$번째 effective transition 진행률, 평형 target ($0\!\le\!\xi\!\le\!1$) \\
$U_j(T),w_j(T)$ & V & 전이 중심 전위, 평형 전이폭 \\
$\mathcal A_j$ & J/mol & 전위 구동력 $=s_{\phi,j}F(V_{n,\drive}-U_j)$ \\
$\chi_j$ & --- & mobility 가속 계수 ($0\le\chi_j\le1$; Ch3 Level B 의 transfer coeff.\ $\beta$ 와 동일물) \\
$\Delta G_{a,j}$ & J/mol & 활성화 자유에너지 $=\Delta H_{a,j}-T\Delta S_{a,j}$ \\
$G$ & J/mol & local activation barrier (mode 별; hidden rate variable) \\
$k_j,\ k_0(T)$ & 1/s & relaxation mobility(속도상수), Eyring prefactor $k_0=(k_BT/h)\kappa$ \\
$L$ & --- & local relaxation length. $q$-좌표서 $L=|I|/(Q_\cell k)$ (무차원) \\
$\rho_G(G)$ & mol/J & activation barrier 분포 \\
$A_L(L)$ & 1/L & relaxation-length spectral density \\
$Q_{j,\tot}$ & C & $j$ transition 의 용량 기여 (방전 기준 $>0$) \\
$Q_\bg(V_n,T)$ & C & 배경 누적 용량; $\partial Q_\bg/\partial V_n$ = 잔류 chemical capacitance \\
$\Theta,\ Q_p$ & --- & 전체 phase progress, 그 용량 scale \\
\bottomrule
\end{longtable}

% ====================================================================
\section{전하 보존식과 내부 전위 $V_n$ (무대; (B) 트렁크)}\label{sec:charge}
% ====================================================================
\textbf{물리.} 음극 내부에 누적된 변환 전하 = 외부로 흘린 전하(보존). 이것이 $V_n$ 을 결정한다
(\AL{9}; \cite{doyle1993,newman}). 배경 저장 $Q_\bg$ 와 staging 진행 $\sum_j Q_{j,\tot}\xi_j$ 의 합:
\begin{equation}
\boxed{\,Q_\cell\,q = Q_\bg(V_n,T) + \sum_{j=1}^{N_p} Q_{j,\tot}\,\xi_j\,.}
\label{eq:charge_balance}
\end{equation}
주어진 $(q,T,\xi_j)$ 에서 식~\eqref{eq:charge_balance} 를 $V_n$ 에 대해 푼다.
\textbf{전체 phase progress 정의(10-pass R9-1)}: staging 기여를 하나로 묶어
$Q_p\Theta\equiv\sum_j Q_{j,\tot}\xi_j$ ($Q_p$ = 그 용량 scale, $\Theta\in[0,1]$ 무차원), 따라서
식~\eqref{eq:charge_balance} 는 $Q_\cell q=Q_\bg+Q_p\Theta$. 이로써 kernel integral 의
$d\Theta_\mathrm{tail}/dq$(\S\ref{sec:kernel})와 charge balance 가 같은 $\Theta$ 로 연결된다.
\textbf{유일·매끄러운 해 조건}(GS-4, 임의 정규화 아님):
\begin{equation}
\frac{\partial Q_\bg}{\partial V_n}\ge \epsilon_Q>0 .
\label{eq:monotone}
\end{equation}
$\xi_j=\xi_{\eq,j}$ 이면 그 해가 평형 곡선 $V_n=V_{n,\OCV}(q,T)$ 이다 — \textbf{외부 OCV lookup 은
이 보존식의 평형 특수해}일 뿐이다(중심식이 OCV 읽기로 회귀하지 않음, P3-2).

% ====================================================================
\section{평형 target $\xi_\eq$ (logistic baseline)}\label{sec:eq}
% ====================================================================
\textbf{보수적 입장.} 평형 target 은 (i) smooth, (ii) $0\le\xi_\eq\le1$, (iii) $V_n$ 단조만
요구한다. grounded baseline 은 graphite intercalation 의 lattice-gas/regular-solution chemical
potential \cite{mckinnon1983,bazant2013} (\AL{3a}): ideal 극한에서
\begin{equation}
\xi_{\eq,j}=\Big[1+\exp\!\big(-(V_n-U_j)/w_j\big)\Big]^{-1},\qquad w_j=RT/F\ (\text{ideal}),
\label{eq:logistic}
\end{equation}
즉 smooth logistic ($0\!\to\!1$ 급점프 아님). $w_j$ 는 nonideality$\cdot$disorder 를 포함하는
effective width 로 둘 수 있어 $RT/F$ 에 강제로 묶지 않는다.
\begin{flagbox}
\textbf{형태는 data 결정(DQ-v3-1).} erf/Gaussian-cumulative isotherm 은 graphite 에 대한
열역학 유도가 문헌에 없어 FLAGGED ansatz(\AL{3b}). 저속 OCV near-peak 감쇠 plot
($\log dQ/dV$ vs $(V-U)$ $\to$ 지수=logistic / vs $(V-U)^2$ $\to$ 가우시안)으로 판별
(\S\ref{sec:falsify}). \textbf{본 장 tail 결론은 평형 형태에 무의존}(post-peak 에서
$d\xi_\eq/dq\simeq0$ 만 사용, \S\ref{sec:rate}). 단, Ch3 detailed balance 의 선형성은
logistic 에 의존하므로 baseline 을 logistic 으로 둔다.
\end{flagbox}

\textbf{저온 거동(서 (1) 의 근거).} 식~\eqref{eq:logistic} 의 peak 형상은
$d\xi_\eq/dV_n=(1/w_j)\xi_\eq(1-\xi_\eq)$ 로 결정되어 peak 높이 $\propto 1/w_j$, 폭 $\propto w_j$.
저온 $\to w_j=RT/F$ 감소 $\to$ \emph{더 좁고 가파른 peak}(꼬리 더 빨리 감쇠). 따라서 평형
열역학은 ``저온 짧은 꼬리''를 예측하며, 관찰(저온 긴 꼬리)은 평형으로 설명 불가 $\Rightarrow$
동역학이 필요하다(\S\ref{sec:rate}).

% ====================================================================
\section{★ 유효 mobility 와 keystone (Ch1 = Level A)}\label{sec:barrier}
% ====================================================================
\textbf{Derivation (활성화, \AL{1}).} transition-state theory \cite{eyring1935}:
$\Delta G_{a,j}(T)=\Delta H_{a,j}-T\Delta S_{a,j}$. 전위 구동력
$\mathcal A_j=s_{\phi,j}F(V_{n,\drive}-U_j)$ 가 활성화 과정에 $\chi_j\mathcal A_j$ 만큼 기여하여
유효 지수를 낮춘다(Br\o{}nsted--Evans--Polanyi/Butler--Volmer; \AL{2}, \cite{evanspolanyi1938,bardfaulkner}):
\begin{equation}
k_j(G,T,V_{n,\drive})=k_0(T)\exp\!\Big[-\frac{G-\chi_j\mathcal A_j}{RT}\Big],\qquad k_0=\tfrac{k_BT}{h}\kappa(T).
\label{eq:rate_const}
\end{equation}

\begin{keybox}
\textbf{★ keystone — Ch1 에서 $\chi_j\mathcal A_j$ 는 ``배리어 낮춤''이 아니라 ``relaxation
mobility 의 전위 의존 \emph{가속}''이다 (Level A).} 그 이유를 학부 수준으로 무비약 전개한다.
2-state 반응속도론에서 net rate 는 $\dot\xi_j=r_j^+(1-\xi_j)-r_j^-\xi_j$. 정상점
$\xi_{\eq,j}=r_j^+/(r_j^++r_j^-)$ 근처 1차 선형화하면 $\dot\xi_j=k_j(\xi_{\eq,j}-\xi_j)$,
$k_j=r_j^++r_j^-$ (\AL{5}). 즉 \textbf{$k_j$ = mobility(평형 접근 속도), $\xi_{\eq,j}$ = target(방향)}.
식~\eqref{eq:rate_const} 처럼 $\chi_j\mathcal A_j$ 가 \emph{$k_j$ 전체}에 들어가면 $r_j^+,r_j^-$ 가
같은 배율로 커져 비 $r_j^+/r_j^-=\xi_\eq/(1-\xi_\eq)$ 가 \emph{불변} — 즉 $\mathcal A_j$ 는
방향성 없는 \textbf{mobility scale}이고, target $\xi_\eq$ 는 열역학(\S\ref{sec:eq})이 전담한다.
\emph{방향성을 가진} barrier lowering(forward 장벽만 $-\beta\mathcal A$, backward 는 $+(1-\beta)\mathcal A$
$\Rightarrow$ 비가 바뀜)은 \textbf{Chapter 3 의 Level B}에서 명시 도입하며, 그때 본 장의 $\chi_j$
는 transfer coefficient $\beta_j$ 와 동일물이 된다. 본 장에서 ``배리어 낮춤''이라는 \emph{방향성
함의} 용어를 쓰지 않는다.
\end{keybox}

\begin{boundbox}
\textbf{Marcus bound(\AL{2}).} 선형 $-\chi_j\mathcal A_j$ 는 작은 구동력의 1차 근사다.
Marcus \cite{marcus1956} 는 곡률을 주므로 $|\mathcal A_j|$ 가 클 때 깨진다; 본 식은 tail 영역
$V_{n,\drive}\approx U_j$ 에서 유효. $G-\chi_j\mathcal A_j<0$ 가 되어도 $\max(\cdot,0)$/clip 으로
자르지 \emph{않는다}(GS-4); 그 영역에서는 transport$\cdot$site$\cdot$diffusion$\cdot$finite-current
같은 \emph{물리적 병목}이 rate 를 제한한다(\S\ref{sec:falsify}; Ch3 전류 보존).
\end{boundbox}
\textbf{물리적 병목의 직렬 합성(10-pass R9-G3; $k_{j,\mathrm{lim}}$ 정의).} 강구동서 활성화가 더 이상
rate-limiting 이 아니면, 유효 mobility 를 \emph{직렬 시간척도 합성}으로 둔다(임의 cap 아님):
\begin{equation}
\frac{1}{k_j}=\frac{1}{k_j^{\mathrm{act}}}+\frac{1}{k_{j,\mathrm{lim}}},\qquad
k_j^{\mathrm{act}}=k_0(T)e^{-(G-\chi_j\mathcal A_j)/RT},
\label{eq:kphys_series}
\end{equation}
$k_{j,\mathrm{lim}}$ 은 transport/site/diffusion/finite-current 병목의 effective rate(독립 물리 근거
필요). $k_{j,\mathrm{lim}}\!\to\!\infty$ 면 $k_j\!\to\!k_j^{\mathrm{act}}$(활성화 지배), $k_j^{\mathrm{act}}\!\gg\!k_{j,\mathrm{lim}}$
면 병목 지배. 이후 본문/Ch6 의 $k_j$ 는 이 $k_j(=k_j^{\mathrm{phys}})$ 를 뜻한다.

% ====================================================================
\section{완화 동역학과 single-mode kernel}\label{sec:rate}
% ====================================================================
한 local mode 의 lag $r=\xi_{\eq,j}-\xi_j$ 는 평형 근처 1차 완화(\AL{5}, \cite{onsager1931,degrootmazur1962}):
\begin{equation}
\frac{d\xi_j}{dt}=k_j(\xi_{\eq,j}-\xi_j),\qquad
\frac{d\xi_j}{dq}=\frac{Q_\cell k_j}{|I|}(\xi_{\eq,j}-\xi_j)\ \ (|I|>0).
\label{eq:relax}
\end{equation}
post-peak 에서 target 이 거의 정지($d\xi_{\eq,j}/dq\simeq0$)하면 1계 선형 ODE 가 풀려
\begin{equation}
r(q)\simeq r(q_a)\exp\!\Big[-\frac{q-q_a}{L}\Big],\qquad \boxed{\,L=\frac{|I|}{Q_\cell\,k_j}\,}
\label{eq:single_kernel}
\end{equation}
이다(차원: $|I|/(Q_\cell k)=$무차원). 이는 관측 tail 전체가 단일 지수라는 뜻이 \emph{아니라}
\textbf{하나의 mode 가 만드는 exponential kernel} 이다. \textbf{저온 거동}: $k_j=k_0e^{-(G-\chi\mathcal A)/RT}$
가 저온서 작아져 $L$ 이 커진다 $\Rightarrow$ 꼬리가 길어진다(관찰 일치, 서 (2)).

% ====================================================================
\section{★ Barrier 분포 $\to$ relaxation-length spectrum ((A) 엔진 이식)}\label{sec:spectrum}
% ====================================================================
\textbf{물리.} 실제 전극은 단일 $G$ 가 아니라 입자크기$\cdot$domain$\cdot$defect 차이로 인한
\emph{barrier 분포} $\rho_G(G)$ 를 가진다(\AL{6}). 관측 tail 의 shape 은 $\rho_G$ 와 같지 않다 —
$G$ 가 먼저 rate 로, 다시 length 로 \emph{지수 변환}되기 때문이다. 식~\eqref{eq:rate_const},
\eqref{eq:single_kernel} 에서
\begin{equation}
\boxed{\,L(G,T,V_{n,\drive})=\frac{|I|}{Q_\cell k_0(T)}\exp\!\Big[\frac{G-\chi_j\mathcal A_j}{RT}\Big]\,.}
\label{eq:L_of_G}
\end{equation}
$G$ 의 작은 분산이 $L$ 에서 \emph{지수적으로 확대}된다 — 이것이 stretched tail 의 근원이다
(긴 꼬리는 $\rho_G$ 가 길어서가 아니라 매핑이 지수이기 때문일 수 있다). 역변환
$G(L)=\chi_j\mathcal A_j+RT\ln[Q_\cell k_0 L/|I|]$, Jacobian $|dG/dL|=RT/L$ 에서 relaxation-length
spectral density:
\begin{equation}
A_L(L;T,V_{n,\drive})=\rho_G\big(G(L)\big)\,\frac{RT}{L}\,A_0(L),
\label{eq:spectrum}
\end{equation}
$A_0$ = per-mode amplitude$\cdot$population weight. ($\int A_L\,dL$ = 총 weight.)
\begin{boundbox}
\textbf{Non-uniqueness(\AL{6}, Plonka \cite{plonka2001}).} stretched kinetics $\to\rho_G$ 역매핑은
\emph{비유일}. 따라서 $\rho_G$ 를 관측 tail 에서 역산하지 않고, 독립 측정한 $\rho_G$ 로부터 tail 을
\emph{forward 예측}만 한다(\S\ref{sec:falsify}, N4).
\end{boundbox}

% ====================================================================
\section{★ 관측 tail = kernel integral}\label{sec:kernel}
% ====================================================================
전체 phase progress 의 post-peak tail 은 local exponential kernel 의 중첩:
\begin{equation}
\boxed{\,\frac{d\Theta_\mathrm{tail}}{dq}\simeq\int_0^\infty A_L(L)\,\frac1L\,
\exp\!\Big[-\frac{q-q_a}{L}\Big]\,dL\,.}
\label{eq:kernel_integral}
\end{equation}
이는 Laplace-type mixture 라 $\rho_G$ 형태에 따라 단일지수$\sim$stretched$\sim$power-law 를 smooth
하게 잇는다. single-mode(식~\eqref{eq:single_kernel}) 는 $A_L=\delta(L-L_0)$ 특수해. 두 시점 kernel
$\mathcal K(q,q')=\int A_L(L)\tfrac1L e^{-(q-q')/L}dL$ 로 쓰면 식~\eqref{eq:kernel_integral} 은
$\mathcal K(q,q_a)$ 이다.

% ====================================================================
\section{Self-consistent closure — Plan B (core) + Plan A (Refs 6/7, 검증 tier)}\label{sec:closure}
% ====================================================================
$\xi_j$ 와 $V_n$ 이 charge balance(식~\eqref{eq:charge_balance})로 결합되어 $\Theta$ 가 적분 안팎에
동시에 나타나는 self-consistent 식이 된다(시간적분형은 인과 상한 $q$ 의 \emph{Volterra} 2종):
\begin{equation}
\Theta(q)=\Theta_0+\int_0^q \mathcal K(q,q')\,[\Theta_\eq(q')-\Theta(q')]\,dq'.
\label{eq:volterra}
\end{equation}

\textbf{Plan B (항상 보존되는 core/validator).} 식~\eqref{eq:volterra} (또는 식~\eqref{eq:relax}
+charge balance)를 \emph{direct $g$-grid 수치적분}으로 푼다. 주어진 $\rho_G$(\AL{6})$\cdot$완화
closure(\AL{5}) 하에서 numerically exact reference 이며, Plan A 의 validator 를 겸한다. 저온
$\to$ large-$L$ $\to$ 긴 꼬리 같은 \emph{물리 결론은 closed-form 없이도 성립}한다.

\textbf{Plan A (사용자 PhD Refs 6/7 \cite{lee2011,son2013}; ratio-substitution closed-form).}
self-consistent 적분을 닫는 \emph{후보}. 방법(사용자 논문 JCP 147,144111 (2017) Eq.(32)$\to$(34)):
미지 $\Theta(q')$ 의 자기참조를 ratio 로 두고, 해석적으로 풀리는 reference $\Theta^\simplecase$
(single-mode baseline)의 비로 치환 — $\Theta(q')/\Theta(q)\approx\Theta^\simplecase(q')/\Theta^\simplecase(q)$:
\begin{equation}
\boxed{\,\Theta_{\mathrm A}(q)=\frac{\Theta_0+\int_0^q \mathcal K\,\Theta_\eq\,dq'}
{1+\dfrac{1}{\Theta^\simplecase(q)}\int_0^q \mathcal K\,\Theta^\simplecase\,dq'}\,.}
\label{eq:closed}
\end{equation}
\begin{boundbox}
\textbf{Refs 6/7 의 물리가정 차이 (그대로 쓰면 안 됨; dossier P3-5(e)).}
(i) Refs 는 \emph{Fredholm}(고정 구간$\cdot$정상상태); 본 식은 \emph{Volterra}(인과). 적용하려면
post-peak 에서 target 동결($d\Theta_\eq/dq\simeq0$) + kernel 국소화로 Fredholm 환원해야 하며, 그
붕괴는 \emph{측정}한다(아래 M1,M5). (ii) ratio ansatz 는 사용자 논문 자신이 ``reaction zone 이
넓으면 부정확''이라 경고 — graphite 에서 \emph{넓은 spectrum = stretched tail(관심 영역)} 이므로
\textbf{closure 가 가장 필요한 곳에서 가장 부정확할 수 있다}. 따라서 Plan B(g-grid)를 reference 로
필수 유지하고 ε 를 운용 $T\cdot$rate 영역에서 측정한다.
\end{boundbox}
\textbf{승격 조건(Plan A $\to$ analytic closure).} M1 fixed-domain Fredholm 재정식 가능 / M2 적분
안팎 동일 function class / M3 ratio 적용 변수 명시 / M4 single-mode 극한서 식~\eqref{eq:single_kernel}
복귀 / M5 정량 오차 $\varepsilon=\lVert\Theta_\mathrm A-\Theta_\mathrm B\rVert/\lVert\Theta_\mathrm B\rVert
\le\varepsilon_\mathrm{tol}$. \textbf{어떤 경우에도 Plan B 가 논리 core 로 보존된다.}

% ====================================================================
\section{ICA mapping 과 측정 축 bridge}\label{sec:ica}
% ====================================================================
$dQ=C_\bg(V_n,T)\,dV_n+Q_p\,d\Theta$ 에서
\begin{equation}
\frac{dQ}{dV_n}=\frac{C_\bg(V_n,T)}{1-Q_p\,(d\Theta/dQ)},\qquad C_\bg\equiv\frac{\partial Q_\bg}{\partial V_n}.
\label{eq:ica}
\end{equation}
측정 축은 $V_{n,\app}=V_n+s_I|I|R_n$ 로 가며(P3-1), voltage-axis tail scale 은
$L_\varphi\simeq|dV_n/dQ|_{q_a}L$. \textbf{$R_n$(전압축 분극)과 $\chi_j$(mobility 가속)의 역할을
혼동하지 않는다}(\S\ref{sec:falsify} 의 co-linearity 주의).

% ====================================================================
\section{★ 피팅 가능한 닫힌 논리식 (N-2 deliverable)}\label{sec:fiteq}
% ====================================================================
본 장의 결과물. inner$\to$outer 평가 순서로 단일 식에 조립한다.
\begin{equation}
\boxed{\;
\frac{dQ}{dV_n}(V_n,T,|I|)=\frac{C_\bg(V_n,T)}{\,1-Q_p\,\dfrac{d\Theta_\mathrm{tail}}{dQ}\,}\;}
\label{eq:fiteq}
\end{equation}
\textbf{평가 순서} (각 단계 본문 식 참조):
\begin{enumerate}[label=\textbf{S\arabic*.},leftmargin=2.2em]
\item charge balance(식~\eqref{eq:charge_balance})로 $V_n(q,T,\{\xi_j\})$ 를 root-find.
\item 평형 target $\xi_{\eq,j}(V_n,T)$ (식~\eqref{eq:logistic}).
\item mobility $k_j=k_0(T)e^{-(G-\chi_j\mathcal A_j)/RT}$, $\mathcal A_j=s_{\phi,j}F(V_{n,\drive}-U_j)$ (식~\eqref{eq:rate_const}).
\item barrier$\to$length $L(G)$ (식~\eqref{eq:L_of_G}) $\to$ spectrum $A_L(L)$ (식~\eqref{eq:spectrum}).
\item tail kernel integral $d\Theta_\mathrm{tail}/dq$ (식~\eqref{eq:kernel_integral}); self-consistency 는
  Plan B(g-grid) 또는 Plan A(식~\eqref{eq:closed}, $\varepsilon$ 통과 시).
\item $dQ/dV_n$ (식~\eqref{eq:fiteq}) $\to$ 측정축은 $dQ/dV_{n,\app}$ (\S\ref{sec:ica}).
\end{enumerate}
\textbf{피팅 파라미터}: $\{U_j,w_j,Q_{j,\tot},Q_\bg\}$(평형, 저속 OCV/GITT), $\{\Delta H_a,\Delta S_a,k_0\}$
(다중 $T$ Arrhenius), $\chi_j$(다중 $|I|$/전위), $\rho_G$(독립 입력 — 역산 금지). \textbf{유효범위}:
$G-\chi_j\mathcal A_j\ge0$(Marcus), $V_{n,\drive}\approx U_j+O(w_j)$. \textbf{P1: solver 코드 구현은 범위 외;
본 식은 논리식으로 제시.}

\textbf{좌표 bridge(10-pass R5-2).} 식~\eqref{eq:fiteq} 의 $d\Theta/dQ$ 는 전체 $\Theta$(peak+tail)의
도함수이며 $d\Theta/dQ=(1/Q_\cell)\,d\Theta/dq$; post-peak 영역에서 $d\Theta/dq\simeq d\Theta_\mathrm{tail}/dq$
(식~\eqref{eq:kernel_integral}). $\rho_G$ 는 독립 입력의 \emph{모수형}(예: $G$ 에 대한 Gaussian/log-normal)
으로 주고 입자분포/EIS/계산으로 제약(역산 금지). $1-Q_p\,d\Theta/dQ>0$ 은 식~\eqref{eq:monotone} 단조성
조건과 동치 영역에서 보장.

\textbf{극한 점검(10-pass R7-1, M4).} 빠른 kinetics $k_j\to\infty$ 면 $L=|I|/(Q_\cell k_j)\to0$,
$A_L\to\delta(L)$, kernel 이 $\delta$ 로 수축해 $d\Theta_\mathrm{tail}/dq\to d\Theta_\eq/dq$ (꼬리 소멸,
평형 추종). 단 이는 bracket 이 $0$ 이 되어서가 아니라 작은 lag $\times$ 큰 $k$ 의 곱이 유한히 남는
asymptotic balance 다(singular perturbation; $d\Theta_\mathrm{tail}/dq$ 가 임의 $\delta$-함수가 되지 않음).
$A_L=\delta(L-L_0)$ 특수해서 single-mode 식~\eqref{eq:single_kernel} 로 정확히 복귀(M4 충족).

% ====================================================================
\section{Tail 의 $T\cdot$전위 의존 = spectrum shift}\label{sec:tailT}
% ====================================================================
식~\eqref{eq:L_of_G} 에서: \textbf{저$T$} $\to k_0$↓ 및 $e^{(G-\chi\mathcal A)/RT}$↑ $\to L$↑ $\to$
$A_L$ 가 large-$L$ 로 이동 $\to$ 긴 꼬리. \textbf{고$T$} $\to L$↓ $\to$ 짧은 꼬리. \textbf{전위
가속}: $\partial\ln L/\partial V_{n,\drive}=-\chi_j s_{\phi,j}F/RT<0$ (forward) $\to$ spectrum 을
short-$L$ 로 이동(state jump 아니라 \emph{spectrum shift}; C-rate 겹침의 근원). 평형폭
$w=RT/F$ 는 저온서 좁아지므로 ``저온 긴 꼬리''는 homogeneous 평형으로 설명 불가, kinetic-spectrum
shift 가 필요함을 \emph{강하게 시사}(확정은 \S\ref{sec:falsify} 통과 전제).
\textbf{Grounding 정정(10-pass R4-2)}: ``고정 barrier 분포 $\to$ \emph{T-의존} stretched 완화''라는
\emph{메커니즘 자체}는 신규가 아니라 disordered/fast-ion-conductor kinetics 에서 확립돼 있다
(연속 relaxation-rate 분포의 합 = stretched \cite{lindsey1980}; 고정 activation-energy 분포가 T-의존
stretching 을 주는 정량 결과 \cite{macdonald2000}; static/dynamic disorder + LFER(=BEP) stretched
kinetics \cite{plonka2001}). 따라서 \AL{8} 은 ``novel''이 아니라 \textbf{BOUNDED — 메커니즘 grounded,
\emph{graphite ICA $dQ/dV$ tail 로의 적용}이 본 작업의 기여}로 표기한다.

% ====================================================================
\section{Falsification 과 비퇴화 discriminator}\label{sec:falsify}
% ====================================================================
관측 tail 을 barrier-spectrum 에 귀속하려면 경쟁원(평형폭$\cdot$분극 $R_n$$\cdot$transport$\cdot$
$\rho_G$ 비유일)을 배제해야 한다(\AL{10}, \cite{dubarry2022,fly2020}).
\textbf{비퇴화 signature}: $L\propto|I|$ 는 transport 도 공유(퇴화). 고유 signature 는
식~\eqref{eq:L_of_G},\eqref{eq:psi_shift_inline} 의 \emph{전위 의존} 항 — tail Arrhenius slope 가
$(V_{n,\drive}-U_j)$ 에 의존. \textbf{가장 약한 지점}: charge-transfer 과전압 $\eta_\mathrm{ct}$ 도
Butler--Volmer 라 전위 지수의존이므로 $\chi_j$ 와 단일 전극서 co-linear 가능 — \textbf{다중
$T\times|I|$ + GITT/pulse} 로 $\eta_\mathrm{ct}$ 를 분리한 \emph{후}에만 $\chi_j$ 귀속 성립.
\textbf{Null 규칙}: (N1) tail의 $\ln$ vs $1/T$ slope 가 전위 무의존이면 transport-dominant 로 기각;
(N2) 전위/$|I|$ 증가에 spectrum shift 없으면 $\chi_j$ 항 부재; (N3) 저속 OCV near-peak 가 저온서
안 좁아지면 disorder 기여 분리; (N4) $\rho_G$ 는 독립 근거로 주고 forward 예측만(역산 금지).
여기서 $\partial\ln L/\partial V_{n,\drive}=-\chi_j s_{\phi,j}F/RT$ (식 재기재 \label{eq:psi_shift_inline}).

% ====================================================================
\section{Chapter 2--6 으로 전달되는 기준식}\label{sec:transmit}
% ====================================================================
\begin{itemize}[leftmargin=1.4em]
\item \textbf{Ch2/Ch4 (발열)}: 직접 입력은 $d\xi_j/dt$(식~\eqref{eq:relax}). Level A $k_j$ 사용 시
  가역열은 평형 온도계수$\cdot$transition entropy(thermo)서, 비가역열은 flux--force $J\mathcal A\ge0$ 서.
  $dV_{n,\app}/dT$ 를 가역열로 직접 쓰지 않음.
\item \textbf{Ch3 (반응속도론)}: 본 장 $k_j$ 는 $r_j^+\!+\!r_j^-$ 의 축약. \textbf{방향성 barrier
  lowering(Level B)} = forward $-\beta_j\mathcal A_j$/backward $+(1-\beta_j)\mathcal A_j$ 는 Ch3 에서
  도입하고, 본 장 $\chi_j=\beta_j$. detailed balance $r_j^+/r_j^-=\xi_\eq/(1-\xi_\eq)$ 와 본 장
  $\xi_\eq$(logistic) 정합. \textbf{검증점}: 본 장 Level A 가 Ch3 Level B 의 정상 target
  $\xi_\mathrm{ss}\to\xi_\eq$ 극한과 일치(P3-6).
\item \textbf{Ch5 (히스테리시스)}: branch 비대칭 = Level B 의 $\beta\cdot$landscape 차이; signed
  current/branch target 은 본 장 $\xi_\eq$(true eq)와 metastable target 을 구분.
\item \textbf{Ch6 (수치/검증)}: Plan B = $g$-grid reference; Plan A = Fredholm/Pad\'e 가속 후보(F1--F5,
  $\varepsilon$); 보유 데이터(반쪽셀 GITT$\cdot$STO-OCV$\cdot$0.1/0.2/0.5/1.0C)로 골격$\to$분극$\to$
  tail 순차 검증.
\end{itemize}

% ====================================================================
\section{자체 검수표 (P3 7항목 + GS + keystone)}\label{sec:selfcheck}
% ====================================================================
\begin{longtable}{p{0.78\textwidth} p{0.14\textwidth}}
\toprule 항목 & 결과 \\ \midrule \endhead
(P3-1) $V_n/V_{n,\app}/V_{n,\drive}/V_{n,\OCV}$ 구분 일관 & 통과 (\S\ref{sec:notation},\ref{sec:ica}) \\
(P3-2) 전하 보존식이 $V_n$ 결정 중심식 (OCV 읽기 회귀 X) & 통과 (\S\ref{sec:charge}) \\
(P3-5) Refs 6/7 = ratio-substitution, 5항목 dossier 기록 + 물리가정 차이 & 통과 (\S\ref{sec:closure}, dossier) \\
(P3-6) Ch1 기준식이 Ch2--5 전달식과 무충돌 (keystone 전파) & 통과 (\S\ref{sec:transmit}) \\
(P3-7) ver.N ↔ Chapter N, $\theta$↔$\xi$ 매핑 명시 & 통과 (\S\ref{sec:notation}) \\
(GS-1) 각 식 prose-먼저 + 무비약 (서, keystone 2-state chain) & 통과 \\
(GS-3) 모든 가정 AL 인용; FLAGGED(erf, novel tail) 미오용 & 통과 \\
(GS-4) clip/cap/softplus/step 정의식 0; Marcus/비유일 유효범위 병기 & 통과 \\
★ keystone: $\chi_j\mathcal A_j$ 를 Ch1 = Level A(mobility 가속)로 확정, barrier lowering = Ch3 Level B & 통과 (\S\ref{sec:barrier}) \\
N-1 단일 서사 / N-2 피팅식 본문 결과 & 통과 (서 / \S\ref{sec:fiteq}) \\
\bottomrule
\end{longtable}

\begin{thebibliography}{99}
\bibitem{eyring1935} H. Eyring, \emph{J. Chem. Phys.} \textbf{3}, 107 (1935).
\bibitem{evanspolanyi1938} M. G. Evans, M. Polanyi, \emph{Trans. Faraday Soc.} \textbf{34}, 11 (1938).
\bibitem{marcus1956} R. A. Marcus, \emph{J. Chem. Phys.} \textbf{24}, 966 (1956).
\bibitem{bardfaulkner} A. J. Bard, L. R. Faulkner, \emph{Electrochemical Methods}, Wiley (Ch.~3).
\bibitem{mckinnon1983} W. R. McKinnon, R. R. Haering, \emph{Mod.\ Aspects Electrochem.} \textbf{15}, 235 (1983).
\bibitem{bazant2013} M. Z. Bazant, \emph{Acc. Chem. Res.} \textbf{46}, 1144 (2013).
\bibitem{onsager1931} L. Onsager, \emph{Phys. Rev.} \textbf{37}, 405 (1931).
\bibitem{degrootmazur1962} S. R. de Groot, P. Mazur, \emph{Non-Equilibrium Thermodynamics}, North-Holland (1962).
\bibitem{plonka2001} A. Plonka, \emph{Dispersive Kinetics}, Springer (2001); A. Plonka, ``Linear free energy relations and reversible stretched exponential kinetics…,'' \emph{J. Phys. Chem. B} \textbf{102}, 5835 (1998).
\bibitem{macdonald2000} J. R. Macdonald, ``Stretched exponentials with $T$-dependent exponents from fixed distributions of energy barriers for relaxation times,'' \emph{Phys. Rev. B} \textbf{61}, 228 (2000).
\bibitem{lindsey1980} C. P. Lindsey, G. D. Patterson, ``Detailed comparison of the Williams--Watts and Cole--Davidson functions,'' \emph{J. Chem. Phys.} \textbf{73}, 3348 (1980). (stretched exp = 연속 relaxation-rate 분포의 합)
\bibitem{lee2011} S. Lee, C. Y. Son, J. Sung, S. Chong, \emph{J. Chem. Phys.} \textbf{134}, 121102 (2011).
\bibitem{son2013} C. Y. Son et al., \emph{J. Chem. Phys.} \textbf{138}, 164123 (2013).
\bibitem{doyle1993} M. Doyle, T. F. Fuller, J. Newman, \emph{J. Electrochem. Soc.} \textbf{140}, 1526 (1993).
\bibitem{newman} J. Newman, K. E. Thomas-Alyea, \emph{Electrochemical Systems}, 3rd ed., Wiley (2004).
\bibitem{dubarry2022} M. Dubarry, D. Anse\'an, \emph{Front. Energy Res.} \textbf{10}, 1023555 (2022).
\bibitem{fly2020} A. Fly, R. Chen, \emph{J. Energy Storage} \textbf{29}, 101329 (2020).
\end{thebibliography}

\end{document}
